\section{Conoscenza di dominio}
\label{sec:domain-knowledge}

\subsection{Ordini di vendita e ordini di produzione}
\label{sec:odv-odp}

Un \textbf{ordine di vendita (ODV)} formalizza la richiesta di fornitura di un cliente e ne definisce le condizioni commerciali. Può comprendere più forniture, articolate in distinte posizioni.

Una \textbf{posizione ODV} è la singola voce dell'ordine di vendita che definisce una specifica fornitura attraverso le specifiche del prodotto, la quantità richiesta e le condizioni di consegna. Rappresenta l'unità di domanda a cui viene associata la produzione ed è identificata dalla coppia numero ODV e numero di posizione.

Un \textbf{ordine di produzione (ODP)} formalizza la richiesta interna di realizzare una determinata quantità di prodotto secondo le specifiche produttive assegnate. Costituisce il riferimento per organizzare e registrare le lavorazioni necessarie alla sua realizzazione.

% Domain sources: StrategicPlannerBeanTech/notes/Email Correspondence
% August 24-26 2026.md; notes/Meeting September 29 2026 Decisions.md.
Nel contesto considerato, ogni ODP è associato a una e una sola posizione ODV, mentre una stessa posizione ODV può essere associata a uno o più ODP. L'ODV esprime quindi la domanda del cliente; gli ODP ne rappresentano la traduzione in richieste di produzione.

% Example verified directly against StrategicPlannerBeanTech/data/raw/2026-09-24/:
% ordini_di_vendita_260915.csv: Numero_OdV=0001113724, Posizione_OdV=000130;
% ordini_di_produzione_260915.csv: ODV=0001113724, POSIZIONE_ODV=000130.
% Join on both identifiers, preserving leading zeros. Both ODP have material
% 3010148, steel 35MNVS5+CR, dimension 91 mm and requested length 7968 mm.
\begin{examplebox}[ex:odv-two-odp]{Una posizione ODV associata a due ODP}

Nell'estrazione ricevuta il 24 settembre 2026, la posizione \texttt{000130} dell'ODV \texttt{0001113724} riguarda il materiale \texttt{3010148}, con data di consegna richiesta pari al 2 gennaio 2026. Alla posizione sono associati i seguenti ordini di produzione:

\begin{center}
\begin{tabular}{ll}
\toprule
\textbf{ODP} & \textbf{Colata} \\
\midrule
\texttt{000010874375} & \texttt{298932} \\
\texttt{000010880354} & \texttt{299558} \\
\bottomrule
\end{tabular}
\end{center}

\end{examplebox}

\subsection{Macrofasi del processo produttivo}
\label{sec:production-phases}

Per realizzare il prodotto richiesto da un ODP, il materiale attraversa un percorso di lavorazioni e controlli, fino al suo approntamento. Le registrazioni che documentano quanto avvenuto durante questo percorso sono dette \textbf{consuntivi di produzione}. Esse associano all'ODP le attività e gli stati registrati, con le relative informazioni temporali e quantitative, e permettono di ricostruirne la storia produttiva.

Una \textbf{macrofase} è un passaggio di questo percorso, descritto a un livello che riassume il dettaglio delle singole operazioni. Può riunire più lavorazioni in un'unica fase, identificare un controllo oppure rappresentare uno stato del materiale, come l'attesa di un esito o il completamento della produzione. Il termine indica quindi il ruolo del passaggio nel processo; non implica che esso corrisponda a una sola macchina.

Nei consuntivi, ogni registrazione riferita a una macrofase ne descrive una specifica occorrenza per un ODP. Una stessa macrofase può comparire più volte nel percorso dello stesso ordine. La successione dei passaggi registrati descrive il \textbf{ciclo produttivo osservato} dell'ODP.

% Inventory checked directly against all records in
% StrategicPlannerBeanTech/data/raw/2026-09-24/Consuntivi_260915.csv,
% field MACHINE_AGGREGATION_CODE: 14 distinct values, 35893 records.
% Counts refer to the final source files before physical-consistency filtering.
% Domain sources:
% - data/raw/2026-09-24/ABS_Flussi_Luna.pptx, slides 1--8;
% - notes/Email Correspondence August 24-26 2026.md;
% - notes/Meeting August 27 2026 Summary.md;
% - notes/Meeting September 29 2026 Decisions.md.
Nei consuntivi dell'estrazione ricevuta il 24 settembre 2026 sono presenti 14 macrofasi distinte. Il loro significato è ricostruito a partire dal documento aziendale \emph{ABS\_Flussi\_Luna} e dalle note di progetto.

\subsubsection{Processo iniziale aggregato}

\begin{description}
\item[\texttt{TRENOLUNA}.] Comprende le lavorazioni in continuo del Treno Luna, dal riscaldo al controllo qualità in linea. Il diagramma aziendale distingue al suo interno la laminazione, il percorso di raffreddamento o trattamento e il controllo del materiale. Nei consuntivi queste attività sono riunite in un'unica macrofase, precedente alle eventuali lavorazioni successive e ai percorsi di recupero del materiale non conforme.
\end{description}

\subsubsection{Lavorazioni e controlli}

\begin{description}
\item[\texttt{BANCHI}.] Indica il controllo del materiale ai banchi, ricondotto nel diagramma aziendale alla voce banco/metalloscopio. Serve a verificare il materiale e a indirizzare gli eventuali interventi di recupero dopo il controllo in linea.

\item[\texttt{TAG}.] Identifica le operazioni di taglio, rappresentate nella documentazione dalla voce \emph{taglio seghetti}. Può comparire nei percorsi di rilavorazione del materiale.

\item[\texttt{MOL}.] Indica la molatura, una lavorazione meccanica del materiale presente nei percorsi di recupero insieme al taglio e alla raddrizzatura.

\item[\texttt{RADDR}.] Indica la raddrizzatura, ossia la lavorazione destinata a correggere la rettilineità del materiale. Nel diagramma compare sia nei percorsi di recupero sia dopo i trattamenti termici.

\item[\texttt{CIRCO}.] È interpretata come il controllo mediante \emph{CIRCOGRAPH}, descritto nel PowerPoint come controllo qualità finale per le rilavorazioni.

\item[\texttt{CONF}.] È interpretata come confezionamento, coerentemente con la voce \emph{Confez} presente nei diagrammi in prossimità dell'uscita del materiale.
\end{description}

Il ruolo di una lavorazione dipende dal percorso del prodotto: la presenza di una macrofase non basta a qualificarla come rilavorazione. La raddrizzatura, per esempio, può essere un passaggio previsto dopo un trattamento termico oppure un intervento di recupero.

\subsubsection{Trattamenti}

Nel PowerPoint i trattamenti termici sono distinti in base alla lavorazione e alla risorsa utilizzata. Il forno a induzione può eseguire il solo rinvenimento o la normalizzazione, oppure il ciclo completo di tempra e rinvenimento: queste alternative impiegano la stessa risorsa. La voce \texttt{DFO} indica invece un trattamento termico esterno a Luna, presente nel percorso di ricottura. Dopo il trattamento, il materiale raggiunge il punto \texttt{HPT} per l'attesa dei risultati del controllo di laboratorio. Nei dati, i passaggi di trattamento sono registrati attraverso le seguenti macrofasi:

\begin{description}
\item[\texttt{TRATT\_LUNA}.] Identifica una macrofase di trattamento riferita a Luna.

\item[\texttt{TRATT\_LFC}.] Identifica un'ulteriore macrofase di trattamento, distinta nei dati da \texttt{TRATT\_LUNA}.
\end{description}

\subsubsection{Controlli di rilascio e stati di avanzamento}

\begin{description}
\item[\texttt{HPT}.] Rappresenta il punto di attesa dell'esito del controllo qualità fuori linea, svolto in laboratorio dopo un trattamento termico. Il rilascio del materiale dipende dal risultato del controllo. Nei consuntivi questa macrofase ha durata registrata nulla: l'attesa dei risultati può comunque separare temporalmente le attività adiacenti.

\item[\texttt{ANTIMIX}.] Nel documento aziendale indica il controllo volto a rilevare mescolamenti qualitativi del materiale. Nelle indicazioni di progetto viene trattata come un punto di controllo o attesa, senza un vincolo di capacità proprio, anche quando il consuntivo riporta una durata positiva. Occorre quindi distinguere il significato fisico del controllo dalla sua rappresentazione nel modello di pianificazione.

\item[\texttt{PRONTO}.] È lo stato terminale che indica l'approntamento del materiale al termine del percorso registrato. Identifica il completamento produttivo rappresentato nei consuntivi, senza attestare di per sé l'avvenuta spedizione al cliente.

\item[\texttt{PRONTO\_QS}.] È una variante dello stato terminale di approntamento. Il significato specifico del suffisso \texttt{QS} non è definito nelle fonti disponibili e viene mantenuto senza attribuirgli una lavorazione aggiuntiva.

\item[\texttt{PRONTO\_QWR}.] È un'ulteriore variante dello stato terminale di approntamento. Anche il significato del suffisso \texttt{QWR} richiede un chiarimento aziendale; nei consuntivi resta distinta dalle altre due macrofasi terminali.
\end{description}

\subsection{Cicli di base e varianti del percorso produttivo}
\label{sec:base-production-cycles}

Il \textbf{ciclo di base} descrive il percorso tecnologico previsto per ottenere un prodotto con le caratteristiche richieste, quando le lavorazioni e i controlli hanno esito regolare. Durante la produzione, il percorso può essere modificato dall'esito dei controlli: una parte del materiale può richiedere verifiche aggiuntive, lavorazioni di recupero o ulteriori trattamenti. Questi passaggi danno origine alle \textbf{varianti} del ciclo.

% Primary source: StrategicPlannerBeanTech/data/raw/2026-09-24/
% ABS_Flussi_Luna.pptx, slide 1 (legend), slides 2--8 (seven flows).
% Markdown sources: notes/Email Correspondence August 24-26 2026.md,
% section "Ciclo base e varianti consuntive"; notes/Meeting August 27
% 2026 Summary.md, sections 5.3 and 6; notes/Meeting September 29 2026
% Decisions.md, section "Sequential process abstraction".
% Table rows describe the main successful path in each slide. Conditional
% recovery paths and the alternative Q+T branch of slide 4 are described below.
La presentazione aziendale \emph{ABS\_Flussi\_Luna} distingue sette flussi, riportati nella \cref{tab:documented-base-cycles}. Per descriverne i passaggi si mantiene la terminologia dei diagrammi: \texttt{LAM} indica la laminazione, \texttt{PLA} la placca di raffreddamento e \texttt{CTR} il controllo in linea. Il forno \texttt{ONA} costituisce un percorso alternativo alla placca di raffreddamento. Le voci \texttt{INDUZIONE R}, \texttt{INDUZIONE N} e \texttt{INDUZIONE Q+T} indicano rispettivamente rinvenimento, normalizzazione e ciclo completo di tempra e rinvenimento. \texttt{DFO} indica un trattamento termico esterno a Luna; \texttt{BORDIONI} identifica gli avvolgitori dell'omonima linea.

\begin{table}[htbp]
\centering
\small
\begin{tabularx}{\linewidth}{@{}p{0.24\linewidth}>{\raggedright\arraybackslash}X@{}}
\toprule
\textbf{Flusso} & \textbf{Percorso principale} \\
\midrule
Naturale &
LAM $\rightarrow$ PLA $\rightarrow$ CTR $\rightarrow$ confezionamento $\rightarrow$ uscita. \\
\addlinespace
ONA &
LAM $\rightarrow$ ONA $\rightarrow$ CTR $\rightarrow$ HPT $\rightarrow$ uscita. \\
\addlinespace
ONA + rinvenimento &
LAM $\rightarrow$ ONA $\rightarrow$ CTR $\rightarrow$ HPT $\rightarrow$ INDUZIONE R $\rightarrow$ HPT $\rightarrow$ ANTIMIX $\rightarrow$ uscita. \\
\addlinespace
ONA + induzione &
LAM $\rightarrow$ ONA $\rightarrow$ CTR $\rightarrow$ HPT $\rightarrow$ INDUZIONE Q+T $\rightarrow$ HPT $\rightarrow$ RADDR $\rightarrow$ CIRCOGRAPH $\rightarrow$ ANTIMIX $\rightarrow$ uscita. \\
\addlinespace
PLA + induzione &
LAM $\rightarrow$ PLA $\rightarrow$ CTR $\rightarrow$ INDUZIONE Q+T o N $\rightarrow$ HPT $\rightarrow$ RADDR $\rightarrow$ CIRCOGRAPH $\rightarrow$ ANTIMIX $\rightarrow$ uscita. \\
\addlinespace
PLA + ricottura DFO &
LAM $\rightarrow$ PLA $\rightarrow$ CTR $\rightarrow$ DFO $\rightarrow$ HPT $\rightarrow$ RADDR $\rightarrow$ CIRCOGRAPH opzionale $\rightarrow$ ANTIMIX $\rightarrow$ uscita. \\
\addlinespace
Bordioni &
LAM $\rightarrow$ BORDIONI $\rightarrow$ confezionamento $\rightarrow$ uscita. \\
\bottomrule
\end{tabularx}
\caption{Flussi descritti nella presentazione aziendale e relativi percorsi principali.}
\label{tab:documented-base-cycles}
\end{table}

\subsubsection{Percorsi principali}

Nel \textbf{flusso naturale}, il materiale attraversa la laminazione, il raffreddamento sulla placca e il controllo in linea. Queste tre fasi sono aggregate nei dati nella macrofase \texttt{TRENOLUNA}. Se il materiale è conforme, prosegue verso il confezionamento e l'approntamento: al livello delle macrofasi, il ciclo di base è quindi \texttt{TRENOLUNA} $\rightarrow$ \texttt{PRONTO}.

Nel \textbf{flusso ONA}, il passaggio nel forno ONA sostituisce il raffreddamento sulla placca. Anche in questo caso, laminazione, passaggio nel forno e controllo in linea sono aggregati nella macrofase \texttt{TRENOLUNA}. Segue il punto \texttt{HPT}, nel quale il materiale attende l'esito del controllo di laboratorio: se favorevole, il materiale raggiunge lo stato \texttt{PRONTO}. Il percorso di base è quindi \texttt{TRENOLUNA} $\rightarrow$ \texttt{HPT} $\rightarrow$ \texttt{PRONTO}.

I quattro flussi con \textbf{trattamenti fuori linea} aggiungono al percorso iniziale il trattamento necessario a ottenere le caratteristiche richieste. Nel percorso ONA con rinvenimento, dopo il primo \texttt{HPT} sono previsti il rinvenimento a induzione, un nuovo \texttt{HPT} e il controllo antimix. Lo stesso diagramma ammette anche un ramo con ciclo completo a induzione, seguito da \texttt{HPT}, raddrizzatura, controllo CIRCOGRAPH e antimix. Il flusso ONA con induzione adotta quest'ultima successione come percorso principale.

Nel percorso PLA con induzione, al raffreddamento e al controllo in linea seguono il trattamento a induzione, completo oppure di normalizzazione, e la successione di attesa dell'esito del controllo di laboratorio, raddrizzatura, CIRCOGRAPH e antimix. Nel percorso PLA con ricottura DFO, il trattamento avviene all'esterno di Luna; al rientro sono previsti \texttt{HPT}, raddrizzatura e antimix, con CIRCOGRAPH indicato come opzionale.

Il \textbf{flusso bordioni} segue una struttura distinta: dopo la laminazione, il materiale passa agli avvolgitori, quindi al confezionamento e all'uscita.

\subsubsection{Diramazioni di controllo e recupero}

Nei diagrammi dei flussi naturale, ONA e con trattamenti fuori linea, un esito non conforme del controllo in linea indirizza il materiale ai banchi di controllo, indicati come \texttt{BNC/MET}. Da questo punto il materiale può essere destinato a rottame, proseguire verso il controllo antimix oppure essere sottoposto a taglio, raddrizzatura o molatura. Queste lavorazioni costituiscono rami alternativi del percorso di recupero. Dopo la molatura è previsto un nuovo passaggio ai banchi, seguito dal controllo antimix. Il materiale recuperato si ricongiunge quindi al percorso successivo previsto per il flusso di appartenenza.

Nel flusso ONA, anche il controllo di laboratorio può determinare una diramazione. Se l'esito è favorevole, il materiale procede verso l'uscita; altrimenti il diagramma prevede un trattamento DFO, un rinvenimento a induzione oppure un ciclo completo a induzione. Seguono un nuovo \texttt{HPT} e il controllo antimix, con la possibilità di un passaggio intermedio di raddrizzatura.

I percorsi possono riguardare soltanto una parte del materiale dello stesso ODP. Nel modello si assume che il materiale già conforme attenda il completamento del percorso della parte da recuperare prima dell'approntamento dell'ordine. Questa ipotesi consente di rappresentare ciascun ODP mediante una sequenza che segue il percorso determinante per il completamento, senza identificare esplicitamente i singoli rami e i loro ricongiungimenti.

\subsubsection{Rappresentazione attraverso le macrofasi}

Il dettaglio dei diagrammi e quello dei consuntivi descrivono lo stesso processo a livelli diversi. Le attività iniziali in continuo sono riunite nella macrofase \texttt{TRENOLUNA}. Per il prodotto laminato grezzo allo stato naturale, le note aziendali individuano quindi il ciclo di base nella successione \texttt{TRENOLUNA} $\rightarrow$ \texttt{PRONTO}. Le verifiche e le lavorazioni aggiuntive compaiono come passaggi intermedi quando si realizzano le corrispondenti varianti.

Il ciclo di base è un riferimento tecnologico. Il ciclo osservato per un ODP documenta invece il percorso registrato per quello specifico ordine, che può comprendere controlli aggiuntivi e recuperi. La frequenza di una variante nei dati non modifica, da sola, la definizione del ciclo di base.

% Example verified against both September raw CSVs:
% Consuntivi_260915.csv and ordini_di_produzione_260915.csv.
% ODP 000010859984: TRENOLUNA > PRONTO;
% ODV=0007001624, POSIZIONE_ODV=000010.
% ODP 000010858580: TRENOLUNA > BANCHI > ANTIMIX > PRONTO;
% ODV=0001117836, POSIZIONE_ODV=000020.
% Both master records: PRODOTTO=LAMINATO BARRE, FINITURA=GREZZO,
% STATO_FISICO=NATURALE. No cause of the extra checks is inferred from rows.
% ODP 000010858580, source CSV lines 22--25: ANTIMIX has
% DURATA_TOT_MIN=18, while PRONTO has DURATA_TOT_MIN=0.
\begin{examplebox}[ex:base-cycle-and-variant]{Macrofasi, ciclo di base e variante nei dati}
Gli ODP \texttt{000010859984} e \texttt{000010858580} appartengono entrambi alla tipologia \emph{laminato barre, grezzo, naturale}. I consuntivi riportano i seguenti percorsi:

\begin{center}
\begin{tabularx}{\linewidth}{@{}l>{\raggedright\arraybackslash}X@{}}
\toprule
\textbf{ODP} & \textbf{Percorso registrato} \\
\midrule
\texttt{000010859984} & TRENOLUNA $\rightarrow$ PRONTO. \\
\addlinespace
\texttt{000010858580} & TRENOLUNA $\rightarrow$ BANCHI $\rightarrow$ ANTIMIX $\rightarrow$ PRONTO. \\
\bottomrule
\end{tabularx}
\end{center}

Il primo ordine presenta la successione del ciclo di base; il secondo include i controlli aggiuntivi ai banchi e all'antimix prima dell'approntamento.

Nel secondo ordine, le macrofasi descrivono quindi il processo iniziale aggregato, il controllo ai banchi, il controllo antimix e lo stato finale di approntamento. Per \texttt{ANTIMIX} sono registrati 18 minuti di durata, mentre \texttt{PRONTO} ha durata nulla.
\end{examplebox}
